\documentclass[10pt,twoside,openright]{book}
\usepackage[papersize={169mm,239mm},top=14mm,bottom=14mm,textwidth=125mm,includehead,includefoot]{geometry}
\usepackage[T1]{fontenc}
\usepackage{newpxtext,newpxmath}
\usepackage{microtype}
\usepackage{booktabs}
\usepackage[hidelinks]{hyperref}
\usepackage{fancyhdr}
\pagestyle{fancy}
\fancyhf{}
\fancyhead[LE,RO]{\thepage}
\fancyhead[RE]{\nouppercase{\leftmark}}
\fancyhead[LO]{\nouppercase{\rightmark}}
\title{A Modular Lund Thesis Starter}
\author{Alex Morgan}
\date{2026}
\begin{document}
\frontmatter
\maketitle
\chapter*{Abstract}
\addcontentsline{toc}{chapter}{Abstract}
This compact starter follows the Lund template's book-sized thesis workflow:
front matter, a kappa-style overview, research chapters, and appendices.
Replace this text with the abstract of your dissertation.
\tableofcontents
\mainmatter
\chapter{Introduction}
\label{ch:introduction}
Long-form research documents are easier to maintain when the front matter,
chapters, papers, and appendices have clear boundaries. The upstream repository
includes a full kappa and papers-based example in \texttt{upstream-example.tex}.
\section{A modular workflow}
Use one file per chapter or included paper, and keep local macros close to the
material that needs them. This makes it practical to reuse a chapter in a
paper-based dissertation.
\chapter{Methods}
\label{ch:methods}
Suppose a review takes $r$ hours per section and a thesis has $n$ sections.
Then the planned effort is
\begin{equation}
  T(n)=T_0+rn,
\end{equation}
where $T_0$ is the fixed setup cost.
\begin{table}[ht]
  \centering
  \caption{A small thesis-planning table.}
  \begin{tabular}{lrr}
    \toprule
    Phase & Sections & Days \\
    \midrule
    Outline & 3 & 1 \\
    Draft & 6 & 4 \\
    Revision & 6 & 2 \\
    \bottomrule
  \end{tabular}
\end{table}
\chapter{Conclusion}
The structure is deliberately small so that it can be extended with chapters,
paper subdirectories, and per-paper bibliographies from the upstream example.
\appendix
\chapter{Submission checklist}
\begin{itemize}
  \item Confirm the degree, department, and defence information.
  \item Check the contents, page numbering, figures, and references.
  \item Apply current Lund University requirements before submission.
\end{itemize}
\backmatter
\end{document}
